\documentclass[10pt,a4paper]{amsart}
\usepackage[margin=19mm]{geometry}
\usepackage{amsmath,amssymb,mathtools}
\usepackage[scr=rsfs,cal=euler]{mathalfa}
\usepackage{xfrac}
\usepackage[unicode,hidelinks,bookmarks=false]{hyperref}
\newcommand{\ie}{i.e.\ }
\newcommand{\cf}{cf.\ }
\newcommand{\resp}{respectively\ }
\newcommand{\Subsection}{\subsection}
\providecommand{\nobreakdash}{\nobreak\hbox{-}}
\setlength{\emergencystretch}{2em}
\multlinegap0pt
\usepackage{bm}
\usepackage{bbm}
\usepackage{dsfont}
\usepackage{xspace}
\usepackage{cancel}
\usepackage{mathtools}
\usepackage{amsthm}
\makeatletter
\@ifundefined{theorem}{\newtheorem{theorem}{Theorem}}{}
\@ifundefined{lemma}{\newtheorem{lemma}[theorem]{Lemma}}{}
\makeatother
\makeatletter
\expandafter\ifx\csname proof\endcsname\relax
\newenvironment{proof}[1][Proof]{\noindent\textit{#1.} }{}
\fi
\makeatother
\title[Spin actions and Polygon spaces]{Spin actions and Polygon spaces}
\author{Eunjeong Lee, Jae-Hyouk Lee}
\date{Source: arXiv:2510.03067v1 (CC BY 4.0)}
\begin{document}
\maketitle

\noindent\emph{Source and licence.} Spin actions and Polygon spaces. Version arXiv:2510.03067v1 (CC BY 4.0), distributed under the Creative Commons Attribution 4.0 International licence, as recorded in the arXiv licence metadata for this exact version and shown on the abstract page https://arxiv.org/abs/2510.03067v1. Full rights notice: licenses/arxiv-2510.03067-CC-BY-4.0.txt.\par\smallskip

\noindent\textit{Editorial scope.} Counted unit: the Lemma whose source label is \texttt{Phik} in the article (source0.tex lines 2433--2443, source label \texttt{Phik}), delivered together with the article's own complete original proof of it (source0.tex lines 2445--2542, one \texttt{proof} environment of 98 lines). No mathematics is written, repaired or re-derived: statement and proof are the article's own text line for line. Required context retained uncounted: surjection (lines 915--943). Every definition, display and label of those lines is retained unchanged. Omitted from this package and not redistributed: 1--914, 944--2432, 2444--2444, 2543--2908. Nothing inside a retained excerpt is omitted or altered: every retained body is delivered exactly as the article's lines are written, so no omission receipt of any kind accompanies this package. Citations: the retained excerpts carry no citation command at all, so no bibliography entry of the article is transcribed and no reference is redistributed. Reference adaptation: none was needed and none was made; every reference inside the delivered unit resolves inside it, so no unresolved reference or citation is produced. Printed slips kept literal: no printed word, symbol, hypothesis or number of the counted statement or of its proof was altered. Layout and class: the article's own class is \texttt{article} with packages including amsfonts, amsmath, amssymb, graphicx; the wrapper is the lane's pinned \texttt{amsart} editorial wrapper at 10pt on A4, which restates the theorem-like environments the retained text needs and adds the article's own macro definitions that the retained text needs (none), with extra wrapper packages (bm, bbm, dsfont, xspace, cancel, mathtools). The delivered numbering is the unit's own. Assets: no retained span contains a figure environment, a graphics-inclusion command, a PDF-page inclusion, a table or a TikZ picture; no image file is copied into this package, the package declares no asset dependency and no image file is redistributed with it. Licence: version arXiv:2510.03067v1 of this article is distributed under the Creative Commons Attribution 4.0 International licence, as recorded in the arXiv licence metadata for this exact version and shown on the abstract page https://arxiv.org/abs/2510.03067v1. A full rights notice is delivered at licenses/arxiv-2510.03067-CC-BY-4.0.txt. Characters: every retained excerpt is pure ASCII, so the delivered engine, XeLaTeX with its default Latin Modern text font, reports no missing character.
\begin{lemma}
\label{surjection} For $\mathbb{F}=\mathbb{R},\mathbb{C},\mathbb{H}$,

(1) The Hopf map $\Phi$ is surjective.

(2) For each $(\lambda,\alpha)\in\mathbb{R}\oplus\mathbb{F}$,
\[%
\begin{cases}
\Phi^{-1}\left(
\begin{pmatrix}
\lambda & \alpha\\
\overline{\alpha} & -\lambda
\end{pmatrix}
\right)  \simeq\mathbb{F}(1),\;\text{for}\;(\lambda,\alpha)\neq(0,0)\\
\Phi^{-1}\left(
\begin{pmatrix}
0 & 0\\
0 & 0
\end{pmatrix}
\right)  =\left\{
\begin{pmatrix}
0\\
0
\end{pmatrix}
\right\}  ,\;\text{for}\;(\lambda,\alpha)=(0,0)
\end{cases}
\]

\end{lemma}

\begin{lemma}
\label{Phik} Let $\mathbb{F}=\mathbb{R},\mathbb{C}$, $\mathbb{H}$ and for each
case, let the corresponding $V$ be $\mathbb{R}^{1+dim_{\mathbb{R}}\mathbb{F}}%
$. Then:

(1) $\Phi^{k}$ is surjective.

(2) For each $[P]\in\mathcal{M}_{k}(V)$, $(\Phi^{k})^{-1}([P])\simeq
\mathbb{F}(1)^{l}$ for some $2\leq l\leq k$. In the case of non-degenerate
$[P]$, we have $l=k$.
\end{lemma}

\begin{proof}
(1) The orthonormality condition of $V_{\mathbb{F}}(2,k)$ corresponds, on the
polygon space side, to the condition that the polygons are closed and have
fixed perimeter $1$. For $A\in V_{\mathbb{F}}(2,k)$, because
\[
I_{2}=AA^{\ast}=%
\begin{pmatrix}
x_{1} & \cdots & x_{k}\\
y_{1} & \cdots & y_{k}%
\end{pmatrix}%
\begin{pmatrix}
x_{1} & \cdots & x_{k}\\
y_{1} & \cdots & y_{k}%
\end{pmatrix}
^{\ast}%
\]
we have
\[
\sum_{i=1}^{k}x_{i}^{2}=\sum_{i=1}^{k}y_{i}^{2}=1,\sum_{i=1}^{k}x_{i}\bar
{y}_{i}=0\text{.}%
\]
Therefore, the polygons%
\[
\Phi^{k}\left(
\begin{pmatrix}
x_{1} & \cdots & x_{k}\\
y_{1} & \cdots & y_{k}%
\end{pmatrix}
\right)  =\left(  \frac{\left\vert x_{1}\right\vert ^{2}-\left\vert
y_{1}\right\vert ^{2}}{2}+x_{1}\bar{y}_{1},\cdots,\frac{\left\vert
x_{k}\right\vert ^{2}-\left\vert y_{k}\right\vert ^{2}}{2}+x_{k}\bar{y}%
_{k}\right)  \in\left(  \mathbb{R}\oplus\mathbb{F}\right)  ^{k}\text{,}%
\]
are closed since
\[
\sum_{i=1}^{k}\left(  \frac{\left\vert x_{i}\right\vert ^{2}-\left\vert
y_{i}\right\vert ^{2}}{2}+x_{i}\bar{y}_{i}\right)  =\frac{1}{2}\left(
\sum_{i=1}^{k}\left\vert x_{i}\right\vert ^{2}-\sum_{i=1}^{k}\left\vert
y_{i}\right\vert ^{2}\right)  +\sum_{i=1}^{k}x_{i}\bar{y}_{i}=0
\]
and have fixed perimeter $1\ $because%
\begin{align*}
\sum_{i=1}^{k}\left\vert \frac{\left\vert x_{i}\right\vert ^{2}-\left\vert
y_{i}\right\vert ^{2}}{2}+x_{i}\bar{y}_{i}\right\vert  &  =\sum_{i=1}^{k}%
\sqrt{\left(  \frac{\left\vert x_{i}\right\vert ^{2}-\left\vert y_{i}%
\right\vert ^{2}}{2}\right)  ^{2}+\left\vert x_{i}\bar{y}_{i}\right\vert ^{2}%
}\\
&  =\sum_{i=1}^{k}\frac{\left\vert x_{i}\right\vert ^{2}+\left\vert
y_{i}\right\vert ^{2}}{2}=1\text{.}%
\end{align*}
Therefore, it is clear that the image of $\Phi^{k}$ is contained in
$\mathcal{M}_{k}(V)$.

Conversely, to show that $\Phi^{k}$ is surjective, we consider the way to
recover an orthonormal $2$-frame from a given polygon in $\mathcal{M}_{k}(V)$.
Each edge of the polygon corresponds to a vector $\vec{v_{i}}$, and finding a
preimage of the polygon amounts to finding a pair $%
\begin{pmatrix}
x_{i}\\
y_{i}%
\end{pmatrix}
\in\mathbb{F}^{2}$ such that $(\pi\circ\Phi)\left(
\begin{pmatrix}
x_{i}\\
y_{i}%
\end{pmatrix}
\right)  =\vec{v_{i}}$. This is exactly the inverse image under $\pi\circ\Phi$
for each $i=1,\cdots,k$, which exists by Lemma \ref{surjection}. Therefore,
$\Phi^{k}$ is surjective onto $\mathcal{M}_{k}(V)$.

The structure of the fiber of $\Phi^{k}$ also follows directly from Lemma
\ref{surjection}. Recall that for $(\lambda,\alpha)\neq(0,0)$, we have
$\Phi^{-1}\left(
\begin{pmatrix}
\lambda & \alpha\\
\overline{\alpha} & -\lambda
\end{pmatrix}
\right)  \simeq\mathbb{F}(1)$, while for $(\lambda,\alpha)=(0,0)$,$\ \Phi
^{-1}\left(
\begin{pmatrix}
0 & 0\\
0 & 0
\end{pmatrix}
\right)  \simeq\left\{
\begin{pmatrix}
0\\
0
\end{pmatrix}
\right\}  $.

However, by the definition of $\mathcal{M}_{k}(V)$, zero polygons (i.e.,
polygons with all edges zero) are excluded. Therefore, none of the entries in
a polygon configuration can correspond to the zero element in $\mathbb{R}%
\oplus\mathbb{F\ }$simultaneously. As a result, for a degenerate polygon
$[P]$, the fiber of $\Phi^{k}$ is isomorphic to $\mathbb{F}(1)^{l}$ for some
$l<k$, while for a non-degenerate polygon, the fiber is isomorphic to
$\mathbb{F}(1)^{k}$.\ \ \ \ \ $\blacksquare$
\end{proof}

\end{document}
